\documentclass[11pt]{article}
\usepackage[T1]{fontenc}
\usepackage{amsmath,amssymb,amsthm}
\usepackage[margin=1in]{geometry}
\usepackage{hyperref}
\newtheorem{theorem}{Theorem}
\newtheorem{lemma}[theorem]{Lemma}
\newtheorem{corollary}[theorem]{Corollary}
\newcommand{\RV}{\mathrm{RV}}
\newcommand{\RK}{\mathrm{RK}}
\newcommand{\ind}{\mathbf 1}
\newcommand{\lean}[1]{\texttt{\detokenize{#1}}}
\title{Random choice versus a common ranking\\in equal-capacity bipartite matching}
\author{MathIsEvenEasier\\\small Prepared with OpenAI Codex (GPT-6 Astra)}
\date{6 October 2026}
\begin{document}
\maketitle

\begin{abstract}
Each job independently chooses a uniform neighborhood of prescribed size
among bins with a common positive integer capacity. We prove that choosing
uniformly among available neighbors gives a stochastically larger matching
than choosing according to one common bin ranking. The comparison holds
for every fixed arrival order and therefore for the random-order model
in Arnosti's equal-capacity conjecture. The proof combines a capped
multinomial association inequality, a two-bin coefficient calculation,
and a coupling indexed by the number of accepted jobs. A Lean formalization
proves the complete finite source-model statement and establishes the
association ingredient directly by induction on the cap.
\end{abstract}

\section{Model and result}\label{sec:model}
There are $m\geq1$ bins, each with capacity $C\geq1$, and $n\geq0$
jobs. Job $j$ has a prescribed degree $d_j\in\{0,\ldots,m\}$.
Independently for each job, its feasible set is a uniform $d_j$-subset of
the bins. First fix an arrival order and a bin-priority order, both
independent of these sets. RANDOM-VERTEX (RV) assigns each arriving job
uniformly among its feasible bins with spare capacity, using fresh
randomness. RANKING (RK) assigns it to the first such bin in the common
priority order. A job with no available feasible bin is rejected.
Let $M_{\RV}$ and $M_{\RK}$ denote the numbers accepted.

\begin{theorem}\label{thm:fullconjecture}
For every choice of $m,C,n$ and degrees as above, and every fixed arrival
and bin-priority order,
\begin{equation}\label{eq:main}
 \Pr(M_{\RV}\geq K)\geq\Pr(M_{\RK}\geq K)
 \qquad(K\in\mathbb N).
\end{equation}
The same inequality holds after independently drawing uniform arrival
and bin-priority permutations, as in Arnosti's model.
\end{theorem}

The probabilities include the random graph and the algorithms' choices.
The source conjecture is the assertion following Theorem~2 of
Arnosti~\cite{arnosti}, for common capacities $C>2$. His graph distribution
is defined by fixed job degrees and independent uniform feasible sets.
His RANDOM-VERTEX implementation uses independent uniform preference
permutations for the jobs: the first member of any nonempty available
set in such a permutation is uniform on that set. Because the current
job's permutation is independent of earlier assignments, this is exactly
the sequential rule used here.

Theorem~\ref{thm:fullconjecture} concerns the equal-capacity random-graph
model. Its assumptions do not cover unequal capacities or comparison
conditional on an arbitrary fixed graph. The proof is in
Sections~\ref{sec:loads}--\ref{sec:coupling};
Section~\ref{sec:formal} describes the machine-checked statement.

\section{A comparison of capped loads}\label{sec:loads}
For vectors of equal total, $x$ majorizes $y$ when every upper partial
sum of the decreasing rearrangement of $x$ is at least the corresponding
sum for $y$. A majorization-increasing statistic therefore increases as
loads become more concentrated.

For $b\geq1$, $q\geq0$ and $0\leq r\leq bq$, write
\[
 \Omega_{b,r,q}=\{x\in\{0,\ldots,q\}^b:\textstyle\sum_i x_i=r\},
 \qquad \mu_{b,r,q}(x)=\frac1{Z_{b,r,q}}\prod_{i=1}^b\frac1{x_i!}.
\]
Here $Z_{b,r,q}>0$ is the normalizing constant. This is the load law of
$r$ independent uniform choices of $b$ bins, conditioned on all loads
being at most $q$: the unconditioned mass at $x$ is
$r!b^{-r}\prod_i1/x_i!$.

\begin{lemma}[Capped association]\label{lem:csassociation}
For symmetric, majorization-increasing $F,G$ on $\Omega_{b,r,q}$,
\[
 \operatorname{Cov}_{\mu_{b,r,q}}(F,G)\geq0.
\]
\end{lemma}
\begin{proof}
Apply Cohen--Sackrowitz~\cite[Lemma 4.1]{cs} to iid variables with
carrier $h_q(k)=1/k!$ on integer $0\leq k\leq q$, zero elsewhere.
The support is an interval and successive positive ratios are
$1/(k+1)$, so the carrier is $PF_2$. Conditioning the total to be $r$
gives exactly $\mu_{b,r,q}$. Their upper sums are the decreasing-order
partial sums defining majorization; Definition~2.2 and Remark~2.5 identify
the relevant increasing functions. The integer-valued extension is stated
on page~809 and justified in Section~5. If needed outside the finite
support, extend a statistic at a partial-sum vector by the maximum of its
values at supported vectors below it, using its minimum value when none
exist. This gives an increasing extension. All expectations are finite.
The degenerate cases are deterministic.
\end{proof}

\begin{lemma}[A shared two-bin calculation]\label{lem:pairmeans}
For $u_j\geq v_j\geq0$, $1\leq j\leq t$, put
\[
 Q(X,Y)=\prod_{j=1}^t(u_jX+v_jY),\qquad
 g_k=\binom tk^{-1}[X^kY^{t-k}]Q\quad(0\leq k\leq t).
\]
The sequence $g_k$ is nondecreasing and discretely convex:
$\Delta g_k\geq0$ and $\Delta^2g_k\geq0$ wherever defined.
\end{lemma}
\begin{proof}
For a uniform permutation $\pi$ of $1,\ldots,t$,
$g_k=\mathbb E[\prod_{\ell\leq k}u_{\pi(\ell)}
\prod_{\ell>k}v_{\pi(\ell)}]$. Direct expansion gives, for $h=1,2$,
\[
 \Delta^h g_k=\mathbb E\left[
 \prod_{\ell\leq k}u_{\pi(\ell)}
 \prod_{k<\ell\leq k+h}(u_{\pi(\ell)}-v_{\pi(\ell)})
 \prod_{\ell>k+h}v_{\pi(\ell)}\right]\geq0.
\]
For $h=2$ the two middle terms have the same expectation by exchanging
positions $k+1,k+2$. This proves both assertions, including zero factors.
\end{proof}

\begin{theorem}[Uniform conditional concentration]\label{thm:allconcentration}
Let $r$ independent categorical choices on $b$ bins have probability
rows nonincreasing in the same bin order. Write $N$ for their load vector,
and suppose $E=\{\max_iN_i\leq q\}$ has positive probability.
For every symmetric majorization-increasing function $F$ on
$\Omega_{b,r,q}$,
\[
 \mathbb E(F(N)\mid E)\geq\mathbb E_{\mu_{b,r,q}}F.
\]
\end{theorem}
\begin{proof}
Let $P(z)=\prod_{j=1}^r\sum_{i=1}^b p_{j,i}z_i$, and let
$\overline P$ be its average over all permutations of the variables.
Symmetrizing the bin labels changes neither $F$ nor the cap event.
Put
\[
 L(x)=\left(\prod_i x_i!\right)[z^x]\overline P(z).
\]
The capped symmetrized load law has density $L/\mathbb E_\mu L$
relative to $\mu=\mu_{b,r,q}$: its mass is $[z^x]\overline P$,
whereas $\mu(x)$ is proportional to $1/\prod_i x_i!$.

We show that $L$ increases in majorization. Consider balancing two
loads $a\geq c+2$ to $a-1,c+1$. In the permutation average, pair
terms that exchange the corresponding variables $X,Y$, and fix the
choices assigned to the other variables. The remaining $t=a+c$
choices contribute $Q(X,Y)+Q(Y,X)$, where
\[
 Q(X,Y)=\prod_{j=1}^t(u_jX+v_jY),\qquad u_j\geq v_j\geq0.
\]
By the common row ordering and Lemma~\ref{lem:pairmeans}, the
normalized coefficients $g_k$ of $Q$ are convex. The contribution
of this pair to $L$ is a nonnegative constant times
\[
 a!c!\,[X^aY^c](Q(X,Y)+Q(Y,X))=t!\,(g_a+g_c).
\]
Convexity gives $g_a+g_c\geq g_{a-1}+g_{c+1}$. The other
coordinate factors are unchanged. Thus balancing decreases $L$;
such transfers generate integer majorization.

Lemma~\ref{lem:csassociation} now yields
\[
 \mathbb E(F(N)\mid E)
   =\frac{\mathbb E_\mu(FL)}{\mathbb E_\mu L}
   \geq\mathbb E_\mu F.
\]
The denominator is positive because $\Pr(E)>0$.
\end{proof}

\begin{lemma}[The uniform next choice minimizes the hazard]\label{lem:nextuniform}
Let independent old choices on $b$ bins have probability rows
nonincreasing in the same order. Condition on every load being at
most $q$, where the conditioning event has positive probability.
Then $p_i:=\Pr(N_i=q\mid\max_jN_j\leq q)$ is nonincreasing in $i$.
Consequently any independent next choice with a nonincreasing row
has saturation probability at least $b^{-1}\sum_i p_i$.
\end{lemma}
\begin{proof}
Take bins $i<j$. Condition on which old choices landed in the pair
and on all other destinations. If there are $t$ choices in the pair,
their probabilities of choosing $i,j$ satisfy $u_\ell\geq v_\ell$.
Thus the count $X$ of choices of $i$ has
$\Pr(X=k)=\binom tk g_k$, with $g$ from Lemma~\ref{lem:pairmeans}.
For $t<q$ neither bin reaches $q$; positive cap probability rules out
$t>2q$. Otherwise $t-q\leq q\leq t$. Lemma~\ref{lem:pairmeans}
and $\binom tq=\binom t{t-q}$ give
\[
 \Pr(X=q)=\binom tq g_q\geq\binom t{t-q}g_{t-q}=\Pr(X=t-q).
\]
Conditioning on the common interval $t-q\leq X\leq q$ preserves
this comparison, as does averaging over the finer information.
Thus $p_i\geq p_j$. For the next row $v$, the two similarly ordered
sequences satisfy
\[
 \sum_i v(i)p_i\geq\frac1b\left(\sum_i v(i)\right)
                                  \left(\sum_i p_i\right)
 =\frac1b\sum_i p_i,
\]
because the difference between the two sides is
$\frac1b\sum_{i<j}(v(i)-v(j))(p_i-p_j)\geq0$.
\end{proof}

For use in the matching process, define the uniform saturation hazard by
\[
 h_{b,C}(r)=\frac1b\,
 \mathbb E_{\mu_{b,r,C-1}}\#\{i:x_i=C-1\}.
\]
It is the expected fraction of bins one choice away from saturation.

\begin{corollary}[The conditional hazard for every capacity]\label{cor:allhazard}
For $b\geq1$, $C\geq1$ and $0\leq r\leq b(C-1)$, take independent
old choices $Z_1,\ldots,Z_r$ and a next choice $Y$ on $b$ bins,
all with probability rows nonincreasing in the same order. If
$N_i=\#\{j:Z_j=i\}$ and $E=\{\max_iN_i<C\}$ has positive probability,
then
\[
 \Pr(N_Y=C-1\mid E)\geq h_{b,C}(r).
\]
Equality holds when every row is uniform. For $C=1$ necessarily
$r=0$ and the hazard is one.
\end{corollary}
\begin{proof}
Put $q=C-1$ and $B(x)=\#\{i:x_i=q\}$. A balancing transfer in
the capped simplex cannot increase $B$, so $B$ is
majorization-increasing. Theorem~\ref{thm:allconcentration} gives
$\mathbb E(B(N)\mid E)\geq\mathbb E_\mu B$.
Lemma~\ref{lem:nextuniform} shows that replacing the independent ordered
next choice by a uniform choice cannot increase its saturation hazard.
That uniform-choice hazard is $\mathbb E(B(N)\mid E)/b$.
By definition $\mathbb E_\mu B/b=h_{b,C}(r)$, proving the claim,
including $r=0$ and $r=bq$.
\end{proof}

\section{Conditioning at an accepted count}\label{sec:history}
Delete degree-zero arrivals and number the remaining arrivals
$1,\ldots,n$. For RANKING fix the priority order $1,\ldots,m$;
relabeling invariance of the neighborhood model makes every fixed
order have the same output law. Let $\sigma_K$ be the arrival accepting
the $K$th job, with $\sigma_K=\infty$ if there is no such arrival.
When $\sigma_K<\infty$, let $A_K$ count the full days just after it.
Put $(\sigma_0,A_0)=(0,0)$. We compare the distribution of this pair
at each $K$; a level never reached is an absorbing cemetery state.
The cemetery state is ordered above every finite state; finite states
are ordered coordinatewise.

\begin{lemma}[The conditional transition at every accepted count]
\label{lem:allhistory}
Fix $K<mC$ and a positive-probability state
$(\sigma_K,A_K)=(s_0,a)$ with $s_0<\infty$.
Put $b=m-a$ and $r=K-aC$.
For either algorithm the next-success index $S=\sigma_{K+1}$ has law
\begin{equation}\label{eq:waitingkernel}
 \Pr(S=s\mid\sigma_K=s_0,A_K=a)=
 \left(\prod_{j=s_0+1}^{s-1} f_j(a)\right)(1-f_s(a)),
 \qquad f_j(a)=\frac{\binom a{d_j}}{\binom m{d_j}},
\end{equation}
for $s_0<s\leq n$, and mass
$\prod_{j=s_0+1}^n f_j(a)$ at infinity.
Conditional also on $S=s<\infty$, the mark increment
$A_{K+1}-a$ is Bernoulli with probability exactly $h_{b,C}(r)$
for RANDOM-VERTEX and at least $h_{b,C}(r)$ for RANKING.
This includes $a=0$ and $b=1$.
\end{lemma}
\begin{proof}
Refine the current state to data $\rho$ specifying the set $D$ of
days full at $\sigma_K$, the
destination of every job assigned to $D$, and the acceptance or
rejection of every arrival through $\sigma_K$. Thus $\rho$ also
determines each full day's filling time and the full-day set
$D_{j-1}$ before every earlier arrival $j$. Let $I$ be the final
$b$ nonfull days and $J$ the accepted jobs whose destinations in
$I$ are not yet specified. There are $|J|=K-aC=r$ such jobs.

For a fixed availability set, let $u_j(i)$ be the unconditional
probability that arrival $j$ chooses available day $i$. If $v$ days
are full, these probabilities are
\[
 u_j(i)=\frac{1-f_j(v)}{m-v}\quad(\RV),\qquad
 u_j(i)=\frac{\binom{m-s_i}{d_j-1}}{\binom m{d_j}}
       \quad(\RK),
\]
where $s_i$ is the position of $i$ among the available days in
priority order. The second formula counts neighborhoods containing
$i$ and avoiding the $s_i-1$ earlier available days; full days
are allowed among the other neighbors. Thus the RV row is uniform
and the RK row is nonincreasing. These formulas depend on which
days are full, not on the partial loads of available days.

For proposed destinations $z=(z_j:j\in J)\in I^J$, write
$N_i(z)=\#\{j\in J:z_j=i\}$. The joint path probability factors as
\begin{equation}\label{eq:refinedfactor}
 \Pr(\rho,z)=\kappa_\rho
      \prod_{j\in J}u_j^{D_{j-1}}(z_j)
      \mathbf1_{\{\max_{i\in I}N_i(z)<C\}},
\end{equation}
where $\kappa_\rho$ contains the probabilities of prescribed
destinations in $D$ and prescribed rejections. Independence of fresh
neighborhoods and tie-breaking gives this product along the prescribed
availability path. Conversely, every $z$ satisfying the cap realizes
that path: loads only increase, so no member of $I$ can have filled
earlier; the assigned job sets in $D$ already fix all other filling
times. The success indicators also fix $\sigma_K$. Thus capped residual
assignments and compatible histories are in bijection; there is no
additional restriction on $z$.

Normalize each restricted row by
\[
 p_{j,i}=\frac{u_j^{D_{j-1}}(i)}
                  {\sum_{\ell\in I}u_j^{D_{j-1}}(\ell)},
 \qquad i\in I.
\]
Every denominator is positive on a positive-probability refinement.
Equation~\eqref{eq:refinedfactor} proves that the residual loads
have the law of independent choices with these rows, conditioned on
their cap. The choices need not remain independent after conditioning.
For RV all rows are uniform on $I$; for RK all decrease in the
same inherited priority order. This argument includes $D=\varnothing$.

Until the next acceptance, rejections leave the loads unchanged.
Each has probability $f_j(a)$, for every residual load vector and
every full-day set of size $a$. This proves
\eqref{eq:waitingkernel}. On each fixed positive refinement, conditional also on $S=s$, the fresh
destination at $s$ is independent of the old residual choices: its row on $I$
is uniform for RV and nonincreasing for RK. The waiting event has
introduced no extra weighting of those loads. Corollary~\ref{cor:allhazard}
therefore gives the asserted mark probabilities on each refinement.
The bound depends only on $K,a,m,C$, so averaging over the refinements
preserves it. With $b=1$ the mark increment is deterministic.
\end{proof}

\section{Coupling the acceptance levels}\label{sec:coupling}
\begin{proof}[Proof of Theorem~\ref{thm:fullconjecture}]
Fix the degree order. At each $K$ we construct a coupling of the two
state distributions, maintaining
\[
 \sigma_{K,\RV}\leq\sigma_{K,\RK},\qquad
 A_{K,\RV}\leq A_{K,\RK}
 \quad\text{whenever }\sigma_{K,\RK}<\infty.
\]
Level zero is deterministic. Suppose the invariant holds at $K<mC$.
If RK has not reached $K$, it cannot reach a higher level; continue
RV with its own conditional law; RK stays in the cemetery state.
Otherwise both current hitting times are finite.

The waiting kernel in Lemma~\ref{lem:allhistory} is increasing in
the starting index $\sigma_K$ and the full-day count $a$: for
$s>\sigma_K$ its cumulative probability is
\[
 \Pr(S\leq s\mid\sigma_K,A_K)
       =1-\prod_{j=\sigma_K+1}^s f_j(a),
\]
and each $f_j(a)$ increases in $a$. Starting earlier adds factors
between zero and one. A shared uniform variable and the two inverse
cumulative distribution functions therefore give
$S_\RV\leq S_\RK$, including the atom at infinity.

If $S_\RK=\infty$, no new mark comparison is required. Otherwise
both next hits are finite. If the old full-day counts differ, each
new count is its old value plus zero or one, so the order persists
under any coupling. If both old counts equal $a$, both current states
have the same $b=m-a$ and $r=K-aC$. Conditional on their respective
next-hit indices, Lemma~\ref{lem:allhistory} gives Bernoulli
increment probabilities $h_{b,C}(r)$ for RV and $p\geq h_{b,C}(r)$
for RK. A fresh shared uniform variable orders the increments.
This covers $a=0$ with no separate initialization, and it covers
the last acceptance at level $mC$ as well.

Each coupled transition has the actual conditional law of the next
state given the current state as its marginal. The law of total
probability therefore gives the correct state distributions at $K+1$.
This argument does not assume that either original state process is
Markov. Induction reaches $mC$, and since $\{M\geq k\}=\{\sigma_k\leq n\}$, the hitting-time order
proves every tail $1\leq k\leq mC$; the other tails are trivial.
Relabeling bins preserves independent uniform neighborhoods and the
number accepted. The fixed-priority comparison therefore holds for
every priority permutation. Averaging the inequality over uniform,
neighborhood-independent arrival and priority permutations proves
the source-model assertion.
\end{proof}


\section{Formal statement and verification}\label{sec:formal}
The finite probabilities in the matching model are rational. The Lean
proof works over $\mathbb Q$ and represents capacity $C$ as $q+1$.
A history is a function
$h:\operatorname{Fin}(n)\to\operatorname{Option}(\operatorname{Fin}(m))$,
where \lean{none} means rejection and \lean{some k} means acceptance
at bin $k$. Before job $j$, bin $k$ is available exactly when its load
is at most $q$. The statistic \lean{acceptedCount} counts the accepted
outcomes, and \lean{matchingTail} sums history mass over
$\{h:\operatorname{acceptedCount}(h)\geq K\}$.

\subsection{Identification of the source law}\label{sec:source-law}
For an original job-labeled graph $G=(G_1,\ldots,G_n)$, define
\[
 U_d(G)=\prod_{j=1}^n
 \frac{\ind\{|G_j|=d_j\}}{\binom m{d_j}}.
\]
For a fixed arrival permutation $\sigma$, fixed ranking $e$, and complete
history $h$, let $A_j(h)$ be its available bins before step $j$. Let
$R_{\RV}(A,S,o)$ be the probability of outcome $o$ under uniform choice
from $A\cap S$, and let $R_e(A,S,o)$ be the indicator that $o$ is the
highest-priority available feasible outcome, including rejection.
The source history laws are explicitly
\begin{align*}
 W_{\RV}(h)
 &=\frac1{n!}\sum_\sigma\sum_G U_d(G)
       \prod_{j=1}^n R_{\RV}(A_j(h),G_{\sigma(j)},h_j),\\
 W_{\RK}(h)
 &=\frac1{n!m!}\sum_\sigma\sum_e\sum_G U_d(G)
       \prod_{j=1}^n R_e(A_j(h),G_{\sigma(j)},h_j).
\end{align*}
The same $e$ occurs in every factor of the RANKING product. Every sum
is finite, including $n=0$ with the empty-product convention. The formal
laws \lean{sourceRVHistory} and \lean{sourceRankingHistory} are these
averages, with permutation counts expressed as finite-type cardinalities.
Both laws are proved nonnegative and normalized.

Reindexing $G$ by $\sigma$ is a bijection and changes its degree vector
to $(d_{\sigma(j)})_j$. For a fixed history, each response then depends
on only one graph coordinate. Expanding the product of sums gives exactly
the previously verified history kernels. This is an algebraic identity;
it does not assert independence of the evolving availability sets.
Relabeling every accepted bin by $e$ preserves loads, availability and
the number accepted, and identifies an arbitrary fixed ranking with the
numerical ranking. These are the two changes of variables used in the
formal source transfer.

\begin{theorem}[Formal source-model comparison]\label{thm:leansource}
For all natural $n,m,q,K$ and every degree function
$d:\operatorname{Fin}(n)\to\mathbb N$ satisfying $d(j)\leq m$,
\[
 \sum_{h:\,\operatorname{acceptedCount}(h)\geq K}W_{\RK}(h)
 \leq
 \sum_{h:\,\operatorname{acceptedCount}(h)\geq K}W_{\RV}(h).
\]
This is \lean{MatchingCapacity.source_matching_tail_comparison}
in \lean{formal/SourceModel.lean}.
\end{theorem}
\begin{proof}
The fixed-degree, fixed-priority theorem gives the comparison of history
kernels. The two bijections above identify those kernels with the literal
graph and priority laws. Finite averaging commutes with the tail sum and
preserves the inequality.
\end{proof}

\subsection{Correspondence with the formal proof}
The main declarations below provide entry points into the certificate.
Their exact file locations and dependencies are listed in
\lean{review/source-model-review.json} and \lean{formal/README.rst}.
\begin{enumerate}
\item Lemma~\ref{lem:csassociation}:
  \lean{cappedAssociation_all} proves the required rational association
  statement directly by induction on the cap. The literature theorem is
  not assumed as an axiom.
\item Lemma~\ref{lem:pairmeans} and Theorem~\ref{thm:allconcentration}:
  \lean{PairCoefficients.lean}, \lean{PairApplications.lean},
  \lean{categoricalDensity_majorization}, and
  \lean{capped_categorical_concentration}.
\item Corollary~\ref{cor:allhazard}:
  \lean{ordered_vs_uniform_saturation}, including the exact uniform
  multinomial reference law.
\item Lemma~\ref{lem:allhistory}:
  \lean{historyFiberEquiv}, \lean{refined_path_factorization},
  \lean{stopping_rv_waiting}, \lean{stopping_ranking_waiting},
  \lean{stopping_rv_saturation}, and \lean{stopping_ranking_saturation}.
\item The coupling in Section~\ref{sec:coupling}:
  \lean{actual_ordered_transition_coupling},
  \lean{stopping_laws_ordered_coupling}, and
  \lean{fixed_priority_matching_tail_comparison}.
\item The source transfer:
  \lean{priority_matching_tail}, \lean{ordered_graph_history_identity},
  \lean{graph_history_kernel_identity}, and
  \lean{source_matching_tail_comparison}.
\end{enumerate}

\subsection{Verification record}
The Azure run of 5 October 2026 used Lean~4.34.0 and Mathlib commit
\nolinkurl{5ed2965256430c3649e86755f9576b54eca72435}.
All 91 positive modules, including the audit module, compiled; the record
contains 714 named theorems and an axiom audit of 929 named declarations.
Only \lean{propext}, \lean{Classical.choice}, and \lean{Quot.sound}
occur as axiom dependencies. The proof uses no \lean{sorry}, custom
axioms, \lean{unsafe}, \lean{native_decide}, or \lean{ofReduceBool}.
The intentionally false \lean{NegativeSource.lean} priority claim was
rejected with an unsolved \lean{False} goal.
The returned source hashes, compiled artifacts, dependency audit and
resource deletion are recorded in
\nolinkurl{formal/verification-source-azure.json}.

This manuscript uses the published association theorem for a shorter
exposition. The formalization supplies a direct proof of the required
finite rational specialization. Its theorem statement and source-model
identification have been checked by the authoring agent against the
primary sources; independent external review has not yet occurred.
The earlier research narrative is preserved in
\nolinkurl{archive/research-note-2026-10-05.tex}.

\begin{thebibliography}{9}
\bibitem{arnosti}
N. Arnosti. Greedy Matching in Bipartite Random Graphs.
\emph{Stochastic Systems} \textbf{12}(2), 133--150, June 2022.
First published online 2 November 2021.
\href{https://doi.org/10.1287/stsy.2021.0082}{doi:10.1287/stsy.2021.0082}.

\bibitem{cs}
A. Cohen and H. B. Sackrowitz. Unbiasedness of Tests for Homogeneity.
\emph{The Annals of Statistics} \textbf{15}(2), 805--816, 1987.
\href{https://doi.org/10.1214/aos/1176350376}{doi:10.1214/aos/1176350376}.
\end{thebibliography}
\end{document}
